\documentclass[10pt,a4paper]{amsart}
\usepackage[margin=19mm]{geometry}
\usepackage{amsmath,amssymb,mathtools}
\usepackage[scr=rsfs,cal=euler]{mathalfa}
\usepackage{xfrac}
\usepackage[unicode,hidelinks,bookmarks=false]{hyperref}
\newcommand{\ie}{i.e.\ }
\newcommand{\cf}{cf.\ }
\newcommand{\resp}{respectively\ }
\newcommand{\Subsection}{\subsection}
\providecommand{\nobreakdash}{\nobreak\hbox{-}}
\setlength{\emergencystretch}{2em}
\multlinegap0pt
\usepackage{mdframed}
\usepackage{mdframed}
\usepackage{mdframed}
\usepackage{mdframed}
\usepackage{amssymb}
\usepackage{tikz}
\usepackage{xcolor}
\usepackage{bm}
\usepackage{mdframed}
\newtheorem{lemma}{Lemma}
\newcommand{\X}{X_\ell}
\title{Kikuchi Graphs of Random Hypergraphs are Approximately Johnson}
\author{Pravesh K. Kothari}
\date{Source: arXiv:2606.08597v1 (CC BY 4.0)}
\begin{document}
\maketitle

\noindent\emph{Source and licence.} Kikuchi Graphs of Random Hypergraphs are Approximately Johnson. Version arXiv:2606.08597v1 (CC BY 4.0), distributed by arXiv under the Creative Commons Attribution 4.0 International licence, http://creativecommons.org/licenses/by/4.0/, as recorded in the arXiv licence metadata for this exact version and shown on the abstract page https://arxiv.org/abs/2606.08597v1. Full licence notice: \texttt{licenses/arxiv-2606.08597-CC-BY-4.0.txt}.\par\smallskip

\noindent\textit{Editorial scope.} Counted unit: the article's lemma lemma (source0.tex lines 559--562). The article's complete original proof of it is delivered with it (source0.tex lines 564--589, one \texttt{proof} environment of 26 lines). No mathematics is written, repaired or re-derived: the statement and the proof are the article's own lines, in the article's own notation, and no retained line is substituted or re-flowed.  No further source-local context is needed: every cross-reference of every retained line resolves inside the two delivered spans.  Nothing was omitted from any retained span, so the package carries no omission record.  No retained line contains a citation, so the wrapper carries no bibliography and no bibliography entry.  No retained line contains a figure environment, an image-inclusion command or drawing code, so no image is delivered and the package declares no asset dependency.  Editorial formatting only, in addition: an AMS \texttt{amsart} preamble that restates the article's own declarations and macros used by the retained lines, namely the environments \texttt{lemma} on separate counters, together with the standard packages those lines need.  The retained environments print in this document as Lemma 1 in order of appearance, whereas the article prints the same items with the numbering of its own full document.  The complete native source of the article as distributed in the arXiv e-print for this exact version is bundled unchanged as \texttt{sources/202342/source0.tex}.
\begin{lemma}
For any $d \leq \ell$, $F_d=\operatorname{span}\{p_R: |R|=d\}$ and
$\dim F_d=\binom nd$. In particular, \(F_\ell=V\).
\end{lemma}

\begin{proof}
We first prove that $\{p_R | |R|=d\}$ are linearly independent on ${[n] \choose \ell}$ by induction on $d$.

For \(d=0\), the claim is clear. Suppose \(d\ge 1\), and assume $\sum_{|R|=d} a_R p_R(S)=0$
for every \(S\in\X\). Fix distinct \(p,q\in[n]\), and let $Y\subset [n]\setminus\{p,q\}$ such that $|Y|=\ell-1$.
Comparing the above relation at \(Y\cup\{p\}\) and at \(Y\cup\{q\}\), we get $\sum_{\substack{U\subset Y\\ |U|=d-1}}
\bigl(a_{U\cup\{p\}}-a_{U\cup\{q\}}\bigr)=0$
for every such \(Y\).

This is a relation among the functions \(p_U\), \(|U|=d-1\), on the \((\ell-1)\)-subsets of the ground set \([n]\setminus\{p,q\}\). Since
$d-1\le \min\{\ell-1,n-\ell-1\}$,
the induction hypothesis gives $a_{U\cup\{p\}}=a_{U\cup\{q\}}$ for every \(U\subset[n]\setminus\{p,q\}\) of size \(d-1\). Hence any two \(d\)-sets differing in exactly one element have the same coefficient. The graph on \(d\)-subsets where two sets are adjacent if they differ in one element is connected, so all \(a_R\) are equal to some constant \(a\). Evaluating the original relation at any \(S\in\X\) gives $0=a\binom{\ell}{d}$.
Thus \(a=0\). This proves linear independence.

Now let \(|Q|<d\). On \(\X\) we have
\[
p_Q(S)
=
\frac{1}{\binom{\ell-|Q|}{d-|Q|}}
\sum_{\substack{R\supset Q\\ |R|=d}} p_R(S).
\]
Indeed, if \(Q\subset S\), then the number of \(d\)-sets \(R\) with $Q\subset R\subset S$
is \(\binom{\ell-|Q|}{d-|Q|}\), while if \(Q\not\subset S\), both sides vanish. Therefore \(F_d\) is spanned by the functions \(p_R\) with \(|R|=d\). Since these functions are linearly independent, $\dim F_d=\binom nd$.
%Finally, \(F_\ell=V\). If \(m=\ell\), then the functions \(p_R\) with \(|R|=\ell\) are the point masses on \(\X\). If \(m=N\), then $\dim F_m=\binom nN=\binom n\ell=\dim V$,
%so again \(F_m=V\).
\end{proof}

\end{document}
